\section{From the compressed defect to the projection trace}
\label{sec:trace-transfer}
The positive contraction used above and the orthogonal theta-image
projection have the same renormalized trace on the original tests.
This requires a smoothed tail estimate, since their unsmoothed ranges
need not be close in operator norm.

Let $V(g)$ denote convolution by $g$ on $L^2(\R,dt)$, and set
$k(t)=\tfrac12e^{-|t|}$ and
$\mathcal G=C_c^\infty(\R)+\C k$. When a multiplicative variable
is used below, $g(y)$ abbreviates $g(\log y)$.

\begin{lemma}\label{lem:smoothed-tail}
For $g\in\mathcal G$, the normalized theta columns satisfy, for every
$B>0$,
\[
 \|V(g)(1-S_L)Z_\Lambda\|_2=O_{\chi,g,B}(L^{-B}).
\]
The estimate holds in the principal two-constraint sector as well
as in every nonprincipal primitive sector.
\end{lemma}
\begin{proof}
Write $\ell(\tau)=L(1/2+i\tau,\chi)$ and $m=\widehat g$.
Convexity in a slightly wider strip and Cauchy's formula give
$|\ell(\tau)|+|\ell'(\tau)|\leq C_\chi(2+|\tau|)^{5/16}$.
For precision, the boundaries $\Re s=-1/16,17/16$ give exponent
$9/16$ and zero respectively. Interpolation bounds the smaller
strip $7/16\leq\Re s\leq9/16$ by exponent $5/16$; a circle of
radius $1/32$ about the critical line gives the derivative bound.
The principal pole is removed before interpolation by a rational
factor bounded above and below on this smaller strip.

For a logarithmic half-line projection $P_+$, the Fourier kernel of
$M_mP_+M_\ell-M_{m\ell}P_+$ is, up to a constant of modulus
$(2\pi)^{-1}$,
\[
 m(\tau)\frac{\ell(\sigma)-\ell(\tau)}{\tau-\sigma}.
\]
The derivative bound controls $|\sigma-\tau|\leq1$. Outside that
strip, the square is bounded by a constant times
$|m(\tau)|^2\bigl((2+|\tau|)^{5/8}|\sigma-\tau|^{-2}
+|\sigma-\tau|^{-11/8}\bigr)$.
Its double integral is finite. Since $m\ell$ is bounded, this
proves that $V(g)P_+E_\chi$ extends boundedly from the raw space.
Translating the half-line conjugates by a unitary phase and leaves
the bound unchanged. The same is true for its complementary
half-line. These assertions apply to $k$, whose multiplier is
$(1+\tau^2)^{-1}$, as well as to a compact smooth $g$.

Put $T_g=V(g)1_{x<\Lambda^{-1}}E_\chi$, initially on its graph
domain. On the complementary raw interval $u>\Lambda$, its kernel,
with input measure $du$ and output measure $dy/y$, is
\[
 K(y,u)=u^{-1/2}\sum_{n>\Lambda u}\chi(n)n^{-1/2}g(ny/u).
\]
For the principal character replace the sum by
\[
 \sum_{n>\Lambda u}n^{-1/2}g(ny/u)
       -\int_{\Lambda u}^\infty v^{-1/2}g(vy/u)\,dv.
\]
The latter integral is precisely the mean subtraction in $E_1$.
Define $F(v)=v^{-1/2}g(v)$ and
$A(s)=|F(s)|+\int_s^\infty|F'(v)|\,dv$.
For a compact smooth $g$, $A$ is bounded near zero and eventually
zero. For $k$, direct substitution gives $A(s)=1$ for $s\leq1$
and $A(s)=s^{-3/2}$ for $s\geq1$. Thus $A\in L^2(ds)$.
Abel summation using bounded character sums, or Euler summation
using the bounded discrepancy of the integers from Lebesgue
measure in the principal case, gives
\[
 |K(y,u)|\leq C_\chi\frac{\sqrt y}{u}A(\Lambda y),
 \qquad
 \|T_g(1-P_\Lambda)\|_2^2\leq C_{\chi,g}\Lambda^{-2}.
\]
All kernel identities hold first on a compact raw core and extend
by the boundedness just proved.

For a selected raw vector let
$W\psi=B_R(P_\Lambda B_RP_\Lambda)^{-1}\psi$.
Then $P_\Lambda W\psi=\psi$ and $\|W\|\leq2$.
In a nonprincipal sector, theta--Poisson gives $T_gW\psi=0$.
In the principal sector the same conclusion follows from
$(W\psi)(0)=\psi(0)=0$. Approximate within the fixed band space
to justify this last identity: evaluation at zero is continuous
there, and the smoothed theta operator is bounded.
Consequently $T_g\psi=-T_g(1-P_\Lambda)W\psi$. Summing over a
raw orthonormal selected basis proves
\[
 \|V(g)(1-S_L)E_\Lambda\|_2^2
       \leq4\|T_g(1-P_\Lambda)\|_2^2
       \leq C_{\chi,g}\Lambda^{-2}.
\]
The upper target tail vanishes because the raw support is bounded
by $\Lambda$, and in the principal sector its mean is zero.
Compose with $G_\Lambda^{-1/2}$ and use
Proposition~\ref{prop:gram-control}, noting $c/\Lambda^2=1/q$.
This proves the assertion with every prescribed logarithmic power.
\end{proof}

\begin{proposition}\label{prop:trace-compression}
For every $h\in C_c^\infty(\R)$,
\[
 \Tr\bigl((Q_\Lambda-S_LQ_\Lambda S_L)V(h)\bigr)\longrightarrow0.
\]
\end{proposition}
\begin{proof}
Write $Z_0=S_LZ_\Lambda$, $Z_1=(1-S_L)Z_\Lambda$.
For $h=g*g^*$, set $A=V(g)^*Z_1$ and $B=V(g)^*Z_0$.
The trace in question is exactly
$\|A\|_2^2+2\Re\Tr(A^*B)$.
Since $Z_\Lambda$ is an isometry,
$\|B\|_2\leq\|V(g)^*S_L\|_2=\sqrt{2L}\|g\|_2$.
The preceding lemma bounds $\|A\|_2$ by $O(L^{-B_0})$ for any
$B_0>0$. Hence the trace tends to zero for $g\in\mathcal G$.
For a general compact smooth $h$, the distributional identity
$(1-\partial_t^2)k=\delta_0$ gives $h=k*(h-h'')$.
Polarization expresses this as a linear combination of four
convolution squares of elements of $\mathcal G$. Apply the
already proved result to each term.
\end{proof}

\subsection{The principal polar coordinates}
The Hilbert space for the full principal image is
$\mathcal H^{\rm full}=L^2(\R,dt)\oplus\C^2$.
For an even raw Schwartz function define
\[
 J\psi=(E_1\psi,\psi(0),I\psi).
\]
Normalized dilation by $a>0$ acts on this space as
$R(a)=V(a)\oplus a^{-1/2}\oplus a^{1/2}$, where
$(V(a)F)(t)=F(t-\log a)$.
Let $P_\Lambda^{\rm full}$ project onto
$J\operatorname{Ran}\Pi_\Lambda$, and identify $Q_\Lambda$
with its extension by zero on the two extra coordinates.

\begin{proposition}\label{prop:polar-trace}
The completion for the norm $\|J\psi\|$ is $\mathcal H^{\rm full}$,
and
\[
 \|P_\Lambda^{\rm full}-Q_\Lambda-P_{\C^2}\|_1\longrightarrow0.
\]
In particular
\[
 \Tr(P_\Lambda^{\rm full}R(h))
 =\Tr(Q_\Lambda V(h))
   +\int_\R h(t)(e^{-t/2}+e^{t/2})\,dt+o(1).
\]
No zero-location hypothesis is used.
\end{proposition}
\begin{proof}
Choose a nonzero smooth raw function supported away from zero
with mean zero. Its centered Mellin transform is entire and
nonzero, and its theta transform is the product of that function
with $\zeta(1/2+i\tau)$, nonzero almost everywhere. Fourier
uniqueness therefore shows that the translates of this theta
image have dense span in $L^2(\R)$. The corresponding raw
dilates still have both moments zero. The moment map
$\psi\mapsto(\psi(0),I\psi)$ is onto $\C^2$: use a bump
away from zero and a bump equal to one near zero. Adding these
two profiles to the dense zero-moment images proves the
completion assertion.

The interior prolate approximation and its two-constraint
correction approximate each fixed smooth zero-moment profile in
theta norm. Hence $Q_\Lambda\to I$ strongly on $L^2$, even
though the selected spaces need not be nested. Apply the same
approximation to fixed profiles $f_1,f_2$ with moment vectors
$e_1,e_2$. Correct their projected moment matrix, which tends
to the identity, to obtain $f_{\Lambda,j}$ with exactly those
moments and $E_1f_{\Lambda,j}\to E_1f_j$. The point-evaluation
estimate used here is the ordinary local band-extension estimate
in the principal approximation, not continuity on a fractional
logarithmic Sobolev space.

The orthogonal complement of the zero-moment image inside the
full image has basis
\[
 w_{\Lambda,j}=((1-Q_\Lambda)E_1f_{\Lambda,j},e_j),\qquad j=1,2.
\]
Strong convergence gives $w_{\Lambda,j}\to(0,e_j)$. Their
$2\times2$ Gram matrix tends to the identity. The formula
$W(W^*W)^{-1}W^*$ for their projection proves operator-norm
convergence to $P_{\C^2}$. Its difference from $P_{\C^2}$ has
rank at most four, so convergence also holds in trace norm.
The operator $R(h)$ is bounded for compactly supported $h$;
pairing with this trace-norm convergence gives the formula.
\end{proof}

\begin{corollary}[A rate for the principal completion]
\label{cor:polar-rate}
For the rational principal cutoff,
\[
 \|P_\Lambda^{\rm full}-Q_\Lambda-P_{\C^2}\|_1
       =O(\Lambda^{-1/2}).
\]
Consequently the error in Proposition~\ref{prop:polar-trace} is
$O_h(\Lambda^{-1/2})$ for every $h\in C_c^\infty(\R)$.
\end{corollary}
\begin{proof}
Choose an even $q\in C_c^\infty((-1,1))$ with $q(0)=0$ and
$Iq=1$. Choose an even Schwartz function $r$ whose Fourier
transform is smooth and supported strictly inside $(-1,1)$,
with $r(0)=1$ and $Ir=0$. For the latter choice, take an even
Fourier bump supported away from zero and normalize its integral.
The two raw profiles
\[
 q_\Lambda(u)=\Lambda^{-1}q(u/\Lambda),\qquad
 r_\Lambda(u)=r(\Lambda u)
\]
have coordinate vectors $(0,1)$ and $(1,0)$, respectively.
Exact dilation covariance gives
\[
 \|E_1q_\Lambda\|_2=\Lambda^{-1/2}\|E_1q\|_2,
 \qquad
 \|E_1r_\Lambda\|_2=\Lambda^{-1/2}\|E_1r\|_2.
\]
Truncate $r_\Lambda$ smoothly at a fixed fraction of the time
radius. Its discarded tail and every required derivative are
smaller than every power of $\Lambda^{-1}$, since the argument
of $r$ there has size comparable to $\Lambda^2$. The untruncated
profile has a fixed relative Fourier margin. Thus the
strict-margin approximation of Lemma~\ref{lem:local-prolate-interior}
projects both profiles into $\operatorname{Ran}\Pi_\Lambda$ with
arbitrarily small power errors in theta norm. The raw mean error
is bounded by $\sqrt\Lambda$ times the raw $L^2$ error; the value
error follows from the ordinary interval derivative estimate
used in that lemma. Both are likewise smaller than every power
of $\Lambda^{-1}$.

Order the projected lifts as $(r_\Lambda,q_\Lambda)$ and invert
their $2\times2$ coordinate matrix, which tends to the identity.
This produces two exact coordinate lifts whose
theta norms are $O(\Lambda^{-1/2})$. Removing their projections
onto the zero-coordinate theta image cannot increase those norms.
The remaining two-dimensional space is therefore the graph of
an operator $X_\Lambda:\C^2\to L^2(\R)$ with
$\|X_\Lambda\|=O(\Lambda^{-1/2})$. Its graph projection differs
from $P_{\C^2}$ in operator norm by $O(\|X_\Lambda\|)$ and has
difference rank at most four. The trace-norm estimate follows,
and pairing with the bounded operator $R(h)$ proves the last
assertion.
\end{proof}
